\chapter{Hoja de Ruta del Aprendizaje}
\label{chap:roadmap}

Dependiendo de tu perfil y objetivos, recomendamos las siguientes rutas de
lectura. Cada ruta omite ciertos capítulos y enfatiza otros.

\section{Ruta A: Científico de Datos}

Esta ruta está diseñada para quienes desean usar LisPy como plataforma de
análisis de datos, computación científica y machine learning.

\begin{center}
\begin{tabular}{ll}
\toprule
\textbf{Parte} & \textbf{Capítulos} \\
\midrule
I -- Fundamentos & Todos \\
II -- El Núcleo Lisp & Capítulos 2.1, 2.2, 2.4 \\
III -- El Puente Python & Todos \\
IV -- Programación Híbrida & Todos \\
VI -- Matemáticas Aplicadas & Todos \\
VIII -- Estudios de Caso & Capítulos 8.3, 8.6 \\
\bottomrule
\end{tabular}
\end{center}

\section{Ruta B: Ingeniero de Software}

Para desarrolladores que buscan explotar la metaprogramación y la ingeniería
de sistemas.

\begin{center}
\begin{tabular}{ll}
\toprule
\textbf{Parte} & \textbf{Capítulos} \\
\midrule
I -- Fundamentos & Todos \\
II -- El Núcleo Lisp & Todos \\
V -- Meta-Programación & Todos \\
VII -- Ingeniería de Sistemas & Todos \\
VIII -- Estudios de Caso & Capítulos 8.2, 8.4, 8.7 \\
\bottomrule
\end{tabular}
\end{center}

\section{Ruta C: Matemático}

Para quienes se enfocan en matemáticas simbólicas, numéricas y sus
aplicaciones.

\begin{center}
\begin{tabular}{ll}
\toprule
\textbf{Parte} & \textbf{Capítulos} \\
\midrule
I -- Fundamentos & Todos \\
III -- El Puente Python & Todos \\
VI -- Matemáticas Aplicadas & Todos \\
VIII -- Estudios de Caso & Capítulo 8.1 \\
\bottomrule
\end{tabular}
\end{center}

\section{Ruta D: Hacker de Lenguajes}

Para entusiastas de la metaprogramación y el diseño de lenguajes.

\begin{center}
\begin{tabular}{ll}
\toprule
\textbf{Parte} & \textbf{Capítulos} \\
\midrule
I -- Fundamentos & Todos \\
II -- El Núcleo Lisp & Todos \\
V -- Meta-Programación & Todos \\
VII -- Ingeniería de Sistemas & Capítulos 7.1, 7.4 \\
\bottomrule
\end{tabular}
\end{center}

\section{Ruta E: Principiante Absoluto}

\begin{enumerate}
\item Completar la Parte I (Fundamentos) completamente.
\item Elegir una ruta especializada según intereses.
\item Volver a capítulos omitidos cuando sea necesario.
\end{enumerate}

Cada ruta está diseñada para ser completa en sí misma, proporcionando tanto
profundidad como amplitud en el área de enfoque.
